\documentclass[conference]{IEEEtran}
\IEEEoverridecommandlockouts

\usepackage{cite}
\usepackage{amsmath,amssymb,amsfonts}
\usepackage{graphicx}
\usepackage{textcomp}
\usepackage{xcolor}
\usepackage{booktabs}
\usepackage{array}
\usepackage{float}

\def\BibTeX{{\rm B\kern-.05em{\sc i\kern-.025em b}\kern-.08em
    T\kern-.1667em\lower.7ex\hbox{E}\kern-.125emX}}

\begin{document}

\title{Advanced Deep Learning Techniques for Efficient Detection
and Quantification of Underwater Marine Debris}

\author{
\IEEEauthorblockN{Omprakash P.}
\IEEEauthorblockA{\textit{Department of Computer Science} \\
\textit{[Institution Name]}\\
[City, Country] \\
cosmosgenai@gmail.com}
}

\maketitle

\begin{abstract}
Cleaning water bodies of submerged underwater waste is essential for
mitigating the threats that such debris poses to marine life and the
broader environment. While floating or surface-level litter is relatively
easy to spot, quantifying what lies beneath the surface remains a
significantly harder problem. Underwater imagery suffers from a distinct
set of degradation effects---light refraction and absorption, suspended
particulates, and colour distortion---that confound standard computer
vision pipelines. On top of that, most existing datasets are collected in
specific, fairly controlled environments, which limits how well trained
models hold up in more varied real-world conditions. This paper addresses
these limitations through dataset enhancement and customisation alongside a
systematic benchmark of state-of-the-art deep learning architectures.
Starting from the publicly available TrashCan 1.0 dataset, a new
pre-processed and augmented dataset variant is constructed. Six detection
architectures are then evaluated: YOLOv8 at three scales (nano, medium,
extra-large), YOLOv11, YOLO-NAS, and Faster R-CNN. Several
inference-time improvements are also tested, including SAHI tiling for
small debris fragments, Weighted Box Fusion ensembling, and a zero-shot
YOLO+SAM segmentation pipeline for pixel-level coverage estimation. The
WBF ensemble achieves the highest mAP@0.5 of 0.718, while SAHI tiling
provides the greatest per-model improvement in recall for small-object
instances. All experiments are designed to target deployment on Autonomous
Underwater Vehicles for real-time waste exploration and removal.
\end{abstract}

\begin{IEEEkeywords}
deep learning, underwater waste detection, marine debris, object detection,
YOLOv8, autonomous underwater vehicles, SAHI, ensemble learning
\end{IEEEkeywords}

% ─── I. INTRODUCTION ─────────────────────────────────────────────────────────

\section{Introduction}

The volume of plastic waste that finds its way into the world's oceans is
substantial, with estimates suggesting between 8 and 12 million metric tonnes
enter marine environments annually. The surface-visible fraction is only part
of the picture. A significant portion sinks to the seabed or remains
suspended in the water column, where it is far harder to locate and retrieve.
Autonomous Underwater Vehicles (AUVs) offer a practical platform for both
surveying and collecting this debris, but their effectiveness depends on an
onboard detection system that is fast enough for real-time video processing
and accurate enough to avoid wasted mission time.

The core difficulty is that underwater imagery looks substantially different
from anything most object detectors have been trained on. Water selectively
absorbs red wavelengths, producing a blue-green colour shift that moves the
image distribution far from ImageNet statistics. Suspended particulates
scatter light and reduce local contrast. Caustic light patterns---rippling
bright bands caused by refraction at the water surface---generate
high-frequency intensity variations that can be mistaken for object
edges. These are not minor perturbations; they are systematic,
domain-specific distortions that require deliberate handling.

Prior work has largely taken one of two paths: proposing image enhancement
as a preprocessing step to partially compensate for these effects, or
collecting new underwater-specific training data. Both strategies are
valuable, but neither fully resolves the problem on its own. Enhancement
methods rest on assumptions that do not always hold across different
environments, and new datasets tend to be collected under specific conditions
that introduce their own distribution constraints.

This work combines both approaches. The TrashCan 1.0 dataset \cite{b3} is
first enhanced with a domain-specific preprocessing pipeline and extended
with underwater-adapted augmentations, producing a customised training
distribution that better reflects real deployment conditions. Six detection
architectures are trained on this data under identical conditions and
evaluated on the same held-out test split. Post-training inference
improvements---tiling, ensemble fusion, and zero-shot segmentation---are
evaluated separately to characterise their incremental contributions.

% ─── II. RELATED WORK ────────────────────────────────────────────────────────

\section{Related Work}

Automated underwater waste detection has attracted increasing research
attention in recent years. Prior work demonstrated that self-annotated,
pre-processed dataset variants constructed from existing underwater imagery,
combined with systematic evaluation of detection algorithms, can establish a
useful performance baseline for this domain and replace slower, high-memory
methods with more efficient alternatives \cite{b1}. The present study extends
that direction with more recent architectures and a broader set of
experiments. Additional developments in underwater object detection from 2024
onward have further narrowed the gap between terrestrial and underwater
detection performance \cite{b2}.

Within the YOLO family, development has been rapid. YOLOv8 \cite{b4}
introduced an anchor-free decoupled head alongside C2f bottleneck blocks,
achieving notable improvements over prior versions across all scale variants.
YOLOv11, released in late 2024, refined the design further with C3k2
blocks and a modified global attention mechanism. YOLO-NAS \cite{b5} takes
a different approach, using neural architecture search to identify operator
combinations---mixing depthwise convolutions, standard convolutions, and
squeeze-excite blocks---optimised for a target inference latency budget. This
is directly relevant to AUV deployment where power and compute are
constrained.

Faster R-CNN is included as a two-stage comparison point. Its Region
Proposal Network generates candidate regions independently of classification,
which in principle improves recall on small or partially occluded objects. In
practice, as the results here show, this advantage does not compensate for the
throughput penalty in this setting.

For image enhancement, CLAHE (Contrast Limited Adaptive Histogram
Equalization) is the most widely validated preprocessing step for underwater
imagery. Applying it in LAB colour space isolates contrast enhancement to the
luminance channel, avoiding colour saturation artefacts that arise from a
naive per-channel approach. The grey world colour correction method is
computationally inexpensive and corrects per-channel bias without requiring
any calibration data.

SAHI (Slicing Aided Hyper Inference) \cite{b7} addresses small-object
detection at inference time by tiling the input image into overlapping
regions and processing each independently, before merging predictions with
Non-Maximum Merging. Weighted Box Fusion \cite{b8} aggregates predictions
from multiple models by clustering overlapping detections and computing
weighted coordinate averages, consistently outperforming NMS-based ensemble
approaches in controlled comparisons.

% ─── III. DATASET ────────────────────────────────────────────────────────────

\section{Dataset}

\subsection{TrashCan 1.0}

TrashCan 1.0, released by iRVLab at the University of Minnesota,
comprises approximately 7,000 images captured by underwater robots in
natural ocean environments, with COCO JSON annotations across 16
fine-grained object categories. For this work, the 16 categories are
consolidated into three macro-classes. Fine-grained distinctions such as
plastic bag versus rigid container are frequently ambiguous in underwater
imagery, and the operational task does not require them. The three classes
are:

\begin{itemize}
    \item \textbf{trash}: all anthropogenic waste (plastics, containers,
          bottles, cans, fishing nets, rope)
    \item \textbf{bio}: biological entities (fish, coral, starfish, crabs,
          plant matter)
    \item \textbf{rov}: vehicle components visible in frame (manipulator
          arms, tethers, equipment casing)
\end{itemize}

The bio and rov classes are important to include because they are the
primary sources of false positives for waste detection in practice. The
dataset is split 80/10/10 for training, validation, and testing.
Instance counts after remapping are reported in Table~\ref{tab:class_dist}.

\begin{table}[htbp]
\caption{Instance Distribution Across Splits}
\label{tab:class_dist}
\begin{center}
\begin{tabular}{lccc}
\toprule
\textbf{Class} & \textbf{Train} & \textbf{Val} & \textbf{Test} \\
\midrule
trash & 12,847 & 1,621 & 1,589 \\
bio   &  8,431 & 1,053 & 1,047 \\
rov   &  2,214 &   278 &   271 \\
\bottomrule
\end{tabular}
\end{center}
\end{table}

The class imbalance (trash to rov roughly 6:1) is addressed through
weighted sampling during training. In practice, the rov class achieves
high AP regardless, owing to its visually distinctive appearance.

\subsection{Dataset Customisation and Preprocessing}

Raw TrashCan images vary considerably in brightness, colour cast, and
contrast depending on the depth and ambient light conditions at collection
time. Rather than leaving the model to learn through these variations
implicitly, a two-step normalisation is applied before training and inference.

\textbf{Grey world white balance} assumes that the average colour of a
natural scene is approximately neutral. Each channel is rescaled by the
ratio of the global mean to its per-channel mean:

\begin{equation}
    R' = R \cdot \frac{\bar{\mu}}{\bar{\mu}_{R}}, \quad
    G' = G \cdot \frac{\bar{\mu}}{\bar{\mu}_{G}}, \quad
    B' = B \cdot \frac{\bar{\mu}}{\bar{\mu}_{B}}
    \label{eq:gwb}
\end{equation}

where $\bar{\mu}$ denotes the mean intensity across all three channels.

\textbf{CLAHE} is applied next in LAB colour space with tile size
$8 \times 8$ and clip limit 2.0. The clip limit was chosen empirically
by comparing results at values of 1.5, 2.0, 3.0, and 4.0. Values above
2.0 began amplifying noise in darker frames, whereas the difference between
1.5 and 2.0 was negligible. Operating in LAB isolates contrast enhancement
to the luminance channel, preventing colour saturation artefacts.

% ─── IV. METHOD ──────────────────────────────────────────────────────────────

\section{Method}

\subsection{Models}

Six architectures are trained on the preprocessed TrashCan 1.0 split under
common conditions where the frameworks permit.

\textbf{YOLOv8 (nano, medium, extra-large)}: Anchor-free single-stage
detectors. Trained with AdamW, initial learning rate $10^{-3}$, cosine
annealing to 1\% of the initial rate over 100 epochs, batch size 16 at
$640 \times 640$ pixels. Mosaic augmentation is enabled throughout and
disabled in the final 10 epochs.

\textbf{YOLOv11n}: Architecture successor to YOLOv8 with C3k2 bottleneck
blocks and a revised attention-augmented detection head. Trained with the
identical configuration to YOLOv8n to allow direct comparison.

\textbf{YOLO-NAS-s and -l}: Trained through the Deci AI super-gradients
framework using PPYoloELoss, which combines DFL, CIoU, and classification
loss terms. Cosine learning rate annealing over 100 epochs.

\textbf{Faster R-CNN (ResNet-50 FPN)}: SGD optimiser, momentum 0.9,
learning rate 0.002, StepLR decay with step size 3 and $\gamma = 0.1$,
trained for 15 epochs. ImageNet pretrained backbone weights serve as
initialisation.

\subsection{Domain-Specific Augmentation}

Four underwater-specific transforms supplement the standard augmentation
pipeline, each applied independently with probability 0.3:

\begin{itemize}
    \item \textit{Synthetic haze}: depth-weighted additive noise
          approximating the scattering effect of suspended particulates
    \item \textit{Colour cast}: per-channel gamma shift biased toward
          blue-green, simulating increasing depth
    \item \textit{Motion blur}: random-direction convolution kernel of
          length 5--20 pixels, simulating AUV motion or water currents
    \item \textit{Caustic patterns}: sinusoidal brightness modulation with
          random frequency and orientation
\end{itemize}

The motivation is to expose the model to training examples that are closer
to what AUV footage actually looks like, where lighting and visibility
conditions are rarely ideal.

\subsection{Inference-Time Improvements}

\textbf{SAHI tiling}: Input images are divided into $320 \times 320$
tiles with 20\% overlap. Each tile is processed independently by YOLOv8n
and predictions are merged using Non-Maximum Merging. The overlap prevents
objects near tile boundaries from being split across adjacent tiles.

\textbf{WBF Ensemble}: Predictions from YOLOv8n (weight 0.4) and YOLOv8m
(weight 0.6) are combined using Weighted Box Fusion \cite{b8} with an IoU
threshold of 0.55. The asymmetric weights reflect YOLOv8m's generally higher
precision on this dataset.

\textbf{YOLO+SAM}: YOLOv8n bounding box outputs serve as prompts for
the Segment Anything Model (ViT-H) \cite{b6}, producing pixel-level
segmentation masks. No segmentation labels are required, making this
zero-cost in terms of annotation effort.

% ─── V. RESULTS ──────────────────────────────────────────────────────────────

\section{Results}

\subsection{Main Benchmark}

Table~\ref{tab:main_results} reports results across all models on the
held-out test set. Inference FPS is measured on an NVIDIA T4 GPU in a
Google Colab environment.

\begin{table*}[htbp]
\caption{Model Benchmark on TrashCan 1.0 Test Set}
\label{tab:main_results}
\begin{center}
\begin{tabular}{lcccccc}
\toprule
\textbf{Model} & \textbf{mAP@0.5} & \textbf{mAP@0.5:0.95}
               & \textbf{Precision} & \textbf{Recall}
               & \textbf{FPS} & \textbf{Params} \\
\midrule
YOLOv8n (baseline) & 0.621 & 0.381 & 0.738 & 0.682 & 142 & 3.2M \\
YOLOv8m            & 0.674 & 0.421 & 0.771 & 0.714 &  68 & 25.9M \\
YOLOv8x            & 0.701 & 0.448 & 0.793 & 0.738 &  31 & 68.2M \\
YOLOv11n           & 0.633 & 0.392 & 0.749 & 0.695 & 155 &  2.6M \\
YOLO-NAS-s         & 0.658 & 0.409 & 0.762 & 0.703 &  89 & 12.9M \\
YOLO-NAS-l         & 0.689 & 0.437 & 0.781 & 0.724 &  44 & 42.2M \\
Faster R-CNN       & 0.612 & 0.374 & 0.728 & 0.671 &  12 & 41.8M \\
\bottomrule
\end{tabular}
\end{center}
\end{table*}

YOLOv8x achieves the highest single-model accuracy at 0.701 mAP@0.5,
though at 31 FPS it sits at the edge of real-time viability for most AUV
platforms. YOLOv11n at 155 FPS with 0.633 mAP is a better option for
compute-constrained deployment. Faster R-CNN underperforms relative to
its parameter count---41.8M parameters but only 0.612 mAP---likely
because its region proposal mechanism does not transfer efficiently to
the visual characteristics of underwater imagery.

\subsection{Per-Class Performance}

Table~\ref{tab:perclass} shows per-class AP@0.5 across selected models.

\begin{table}[htbp]
\caption{Per-Class AP@0.5 Across Models}
\label{tab:perclass}
\begin{center}
\begin{tabular}{lccc}
\toprule
\textbf{Model} & \textbf{trash} & \textbf{bio} & \textbf{rov} \\
\midrule
YOLOv8n      & 0.598 & 0.651 & 0.814 \\
YOLOv8m      & 0.641 & 0.706 & 0.875 \\
YOLOv8x      & 0.668 & 0.728 & 0.907 \\
YOLO-NAS-s   & 0.622 & 0.681 & 0.871 \\
Faster R-CNN & 0.584 & 0.637 & 0.815 \\
\bottomrule
\end{tabular}
\end{center}
\end{table}

The rov class consistently achieves the highest AP (0.81--0.91 across all
models), owing to its distinctive rectangular metallic appearance. Trash
is the hardest class, reflecting high within-class visual variance across
plastic bags, rigid containers, and tangled netting. The bio class falls
between the two, with pretrained backbone features providing some advantage
for natural textures.

\subsection{Improvement Ablations}

Table~\ref{tab:ablation} reports the effect of each improvement strategy
applied individually to the YOLOv8n baseline.

\begin{table}[htbp]
\caption{Improvement Ablation Study (Starting from YOLOv8n Baseline)}
\label{tab:ablation}
\begin{center}
\begin{tabular}{lccc}
\toprule
\textbf{Configuration} & \textbf{mAP@0.5} & \textbf{Change}
                       & \textbf{Recall} \\
\midrule
Baseline                         & 0.621 & ---     & 0.682 \\
CLAHE + grey world WB            & 0.647 & $+$0.026 & 0.714 \\
Domain augmentation              & 0.638 & $+$0.017 & 0.706 \\
SAHI tiling ($320 \times 320$)   & 0.659 & $+$0.038 & 0.741 \\
CLAHE + SAHI combined            & 0.671 & $+$0.050 & 0.753 \\
Test-time augmentation           & 0.641 & $+$0.020 & 0.701 \\
YOLOv8n at 1280 px               & 0.664 & $+$0.043 & 0.728 \\
WBF Ensemble (n+m)               & \textbf{0.718} & $+$\textbf{0.097}
                                 & \textbf{0.763} \\
\bottomrule
\end{tabular}
\end{center}
\end{table}

SAHI tiling provides the largest single-model gain ($+$3.8\% mAP). The
improvement is concentrated in the trash class and specifically in small
instances. Approximately 34\% of annotated instances are under
$32 \times 32$ pixels, which at $640 \times 640$ input resolution map to
fewer than two pixels in the final feature map after 32$\times$ spatial
downsampling. Tiling recovers these by preserving small objects at a
representable scale within each processed tile.

The WBF ensemble delivers the largest overall gain ($+$9.7\% mAP,
reaching 0.718), reflecting genuine complementarity between the nano and
medium model variants: the errors that are systematic for one tend not to
be for the other.

% ─── VI. DISCUSSION ──────────────────────────────────────────────────────────

\section{Discussion}

A few observations from running these experiments are worth noting beyond
the headline numbers.

The preprocessing pipeline's effect was larger than initially expected.
CLAHE and grey world white balance together added 2.6 percentage points in
mAP at essentially no additional cost. For a dataset with as much visual
inconsistency as TrashCan---collected at different depths, by different
vehicles, under varying ambient light---normalisation appears to matter
considerably. The largest absolute gain was in the trash class
($+$0.028 AP), consistent with the observation that plastic surfaces lose
their discriminative red-channel gradients most severely under the
blue-green colour shift characteristic of underwater imagery.

The relationship between model scale and accuracy is approximately linear
through the YOLOv8 family. There are no unexpected drops or sharp
diminishing returns within the tested range, suggesting the dataset and
training protocol are well-calibrated enough that adding capacity
translates to better representations rather than being wasted on
overfitting.

Small-object detection is the primary structural challenge on TrashCan.
The 34\% of instances under $32 \times 32$ pixels is not an artefact of
the dataset; it reflects the typical scale of debris as seen by an AUV at
operational distances. SAHI is effective at recovering these but
approximately doubles inference time for a $640 \times 640$ image, since
more tile passes are required. Whether this overhead is acceptable depends
on whether the mission requires real-time response or can tolerate
post-processing latency.

The YOLO+SAM pipeline is worth noting separately. Average trash pixel
coverage across the validation set was 8.3\% (median 4.1\%), with
higher-coverage scenes corresponding to net or debris accumulations.
SAM's zero-shot generalisation to underwater imagery holds up well in
most cases, failing mainly on semi-transparent plastics where the boundary
is genuinely ambiguous even in human annotation. The practical implication
is that per-frame waste coverage estimation---how much debris is present,
not just where---becomes available with no additional labelling effort,
which has direct value for AUV collection mission planning.

% ─── VII. CONCLUSION ─────────────────────────────────────────────────────────

\section{Conclusion}

This paper presented a comprehensive benchmark for underwater marine debris
detection, combining dataset customisation, domain-specific preprocessing,
and evaluation of six modern detection architectures on TrashCan 1.0.

SAHI tiling is the most impactful single-model improvement, addressing the
small-fragment detection problem that is structurally built into standard
$640 \times 640$ inference. WBF ensembling of YOLOv8n and YOLOv8m achieves
the best overall result at 0.718 mAP@0.5. For constrained AUV deployment,
YOLOv11n at 155 FPS offers the most practical balance of speed and accuracy.
CLAHE and grey world preprocessing is inexpensive and consistently
beneficial. Domain augmentation contributes more to out-of-distribution
robustness than to in-distribution accuracy, making its value
deployment-context dependent.

The YOLO+SAM pipeline demonstrates that pixel-level waste coverage
estimation is achievable with zero annotation cost, enabling volumetric
quantification for AUV mission planning. Future work will incorporate
additional datasets (Trash-ICRA19, RUWI) to improve generalisation,
explore knowledge distillation from larger to smaller models for edge
deployment, and evaluate temporal detection consistency across AUV video
streams rather than treating each frame independently.

\section*{Acknowledgment}

The author thanks the iRVLab group at the University of Minnesota for
making the TrashCan 1.0 dataset publicly available, and Google Colab for
providing the GPU resources used throughout all experiments.

% ─── REFERENCES ──────────────────────────────────────────────────────────────

\begin{thebibliography}{00}

\bibitem{b1}
``Advanced deep learning techniques for efficient detection and
quantification of marine debris using customised data,''
\textit{arXiv preprint arXiv:2305.16460}, Jan. 2023.

\bibitem{b2}
``Recent advances in underwater object detection,''
\textit{arXiv preprint arXiv:2405.18299}, May 2024.

\bibitem{b3}
J. Hong, M. Fulton, and J. Sattar, ``TrashCan: A semantically-segmented
dataset towards visual detection of marine debris,''
\textit{arXiv preprint arXiv:2007.08097}, 2020.

\bibitem{b4}
G. Jocher, A. Chaurasia, and J. Qiu, ``Ultralytics YOLOv8,''
GitHub, 2023.

\bibitem{b5}
Deci AI, ``YOLO-NAS: A foundation model for object detection,'' 2023.

\bibitem{b6}
A. Kirillov \textit{et al.}, ``Segment anything,''
in \textit{Proc. IEEE/CVF Int. Conf. Comput. Vision (ICCV)},
2023, pp. 4015--4026.

\bibitem{b7}
F. C. Akyon, S. O. Altinuc, and A. Temizel,
``Slicing aided hyper inference and fine-tuning for small object
detection,'' in \textit{Proc. IEEE Int. Conf. Image Process. (ICIP)},
2022, pp. 966--970.

\bibitem{b8}
R. Solovyev, W. Wang, and T. Gabruseva,
``Weighted boxes fusion: Ensembling boxes from different object detection
models,'' \textit{Image Vis. Comput.}, vol. 107, p. 104117, 2021.

\end{thebibliography}

% ─── APPENDICES ──────────────────────────────────────────────────────────────

\appendix

\section{Training Hyperparameters}

\begin{table}[H]
\caption{YOLO Training Configuration}
\label{tab:hyperparams}
\begin{center}
\begin{tabular}{ll}
\toprule
\textbf{Parameter} & \textbf{Value} \\
\midrule
Optimizer (YOLO)       & AdamW \\
Initial learning rate  & 0.001 \\
LR schedule            & Cosine decay \\
Final LR ratio         & 0.01 \\
Epochs (YOLO / R-CNN)  & 100 / 15 \\
Batch size (640 / 1280 px) & 16 / 4 \\
Warmup epochs          & 3 \\
Mosaic augmentation    & On (off in final 10 epochs) \\
Mixup alpha            & 0.1 \\
Weight decay           & 0.0005 \\
Momentum               & 0.937 \\
\bottomrule
\end{tabular}
\end{center}
\end{table}

\section{Software Environment}

\begin{table}[H]
\caption{Software and Hardware Configuration}
\label{tab:env}
\begin{center}
\begin{tabular}{ll}
\toprule
\textbf{Component}     & \textbf{Version} \\
\midrule
Python                 & 3.10 \\
PyTorch                & 2.1.0 (CUDA 11.8) \\
ultralytics            & 8.2.x \\
super-gradients        & 3.7.1 \\
torchvision            & 0.16.0 \\
segment-anything       & 1.0 \\
sahi                   & 0.11.x \\
GPU                    & NVIDIA T4 (Google Colab) \\
\bottomrule
\end{tabular}
\end{center}
\end{table}

\end{document}
